% document formatting
\documentclass[10pt]{article}
\usepackage[utf8]{inputenc}
\usepackage[left=1in,right=1in,top=1in,bottom=1in]{geometry}
\usepackage[T1]{fontenc}
\usepackage{xcolor}

% math symbols, etc.
\usepackage{amsmath, amsfonts, amssymb, amsthm}

% lists
\usepackage{enumerate}

% images
\usepackage{graphicx} % for images
\usepackage{tikz}

% code blocks
% \usepackage{minted, listings} 
\usepackage[colorlinks=true, urlcolor=blue]{hyperref}


% verbatim greek
\usepackage{alphabeta}

\graphicspath{{./assets/images/Week 2}}

\title{02-719 Week 2 \\ \large{Genomics and Epigenetics of the Brain}}
 
\author{Aidan Jan}

\date{\today}

\begin{document}
\maketitle

\section*{Basic Experiment Setup}
Suppose we are running an experiment with 3 different animals, in two conditions.  Let $A$, $B$, $C$ be in condition A and $D$, $E$, $F$ be in condition B.  Let a scale from 0 to 10 define the amount of open chromatin.  Suppose we run our experiment, and we get:
\begin{center}
\begin{tabular}{c|c}
Species x Condition & Open Chromatin \\ \hline
$A$ & 2 \\
$B$ & 4 \\
$C$ & 3 \\
$D$ & 7 \\
$E$ & 10 \\
$F$ & 8 \\
\end{tabular}
\end{center}
From this, we can then compare the effects on the two different conditions on epigenetics.  Some statistics we can use to quantify the difference are
\begin{itemize}
	\item Mean and Std. Dev
	\item t-test
	\item Linear Regression
	\item etc.
\end{itemize}
Now, if today we want to understand the conditions' effects on multiple genes, we can create matrix where there are $g$ rows, where $g$ is the number of genes we care about, and $c$ number of columns, where $c$ is the number of species x conditions.  Data can then be compared and processed efficiently using a computer program.

\section*{Introduction to R Programming Language}
\begin{itemize}
	\item Very similar to Matlab
	\item Handles matrix operations very well
	\begin{itemize}
	    \item For loops are quite slow
    \end{itemize}
    \item R is open source
    \begin{itemize}
	    \item The underlying code is available
	    \item No license required
    \end{itemize}
    \item Many people have put together statistical and genomic packages in for R
    \begin{itemize}
	    \item \href{http://bioconductor.org/}{Bioconductor} - Open source software for bioinformatics
	    \item Limma - linear models for microarray (and RNA-Seq)
    \end{itemize}
\end{itemize}

\subsection*{Data / Class Types in R}
\begin{itemize}
	\item Boolean, Integer, Numeric (double), String, Factor
	\item Vector - 1D, every element same data type
	\item List - 1D, every element can be different type
	\item Matrix - 2D, every element same data type
	\item DataFrame - 2D, every column same data type
	\item Factor - Used for representing qualitative variables in statistical operations
\end{itemize}

\subsection*{Assigning Different Types}
\begin{itemize}
	\item \texttt{x <- 1}: assignment operator
	\item \texttt{class(x)}: identify data type with class function
	\item \texttt{x <- c(1, 2, 3)}: assign a vector
	\item \texttt{x <- matrix(c(1, 2, 3, 4, 5, 6), 2, 3)}: assign a matrix.  This example makes a matrix of 3 rows, 2 cols.  Numbers fill rows first (1, 2, 3) in col 1.
	\item \texttt{f <- data.frame(num=c(1, 2, 3, 4), letter=c("a", "b", "c", "d"));}: sets data in a data frame.  column name is the key.
	\item \texttt{rownames(f) <- c("first", "second", "third", "fourth");}: sets row names for a dataframe
\end{itemize}

\section*{Gene Ontology}
\begin{itemize}
	\item When we perform a ChIP-Seq experimnet, have peaks (signal) that can potentailly cover any part of the entire genome, not just genes.
	\item The GREAT gene ontology analysis tool is designed to handle this case.
\end{itemize}

\subsection*{Expected Counts}
\textbf{Problem:}  We have a genome that contains 100 genes.
\begin{itemize}
	\item 10 of those genes are involved in DNA replication
	\item We conduct an experiment where key replication gene is knicked out
	\item As a result, 20 genes change in expression
	\item 8 of those genes are involved in replication
	\item What is the significance?
\end{itemize}
To solve:  Let $n$ be number of genes (20), $k$ be number of genes in category (8), and $p$ be probability that any given gene is in the category:  (10/100 = 0.1).\\\\
By simple probability, we would expect that (20 genes $\times$ 0.1 probability = ) 2 genes are in the category.  What we observe was 8 genes, which is much higher than the predicted 2 genes.  $P_{\text{binom}} = 0.0004$.

\subsection*{Binomial P-value}
\begin{itemize}
	\item We can do statistical enrichments based on groups of genes:
	\begin{itemize}
	    \item How many genes do I see in a category?
	    \item How many genes would I expect by chance?
    \end{itemize}
    \[\text{Pr}(k;n,p) = {n \choose k} \cdot p^k \cdot (1 - p)^{n - k}\]
    \item In the formula above,
    \begin{itemize}
	    \item $n$ = number of trials  (number of genes)
	    \item $k$ = number of successes  (number of genes in category)
	    \item $p$ = probability of success  (probability that any given gene is in the category)
    \end{itemize}
\end{itemize}

\subsection*{Mapping Regulatory Elements to Genes}
The assumptions of the binomial test may not be valid.
\begin{center} 
	\includegraphics*[width=\textwidth]{W2_1.png} 
\end{center}

\subsection*{GREAT}
Solution: Partition the genome based on the gene ontology category of the nearest gene.  Then, adjust the probability of a hit based on the proportion of the genome that is covered, not the total number of genes in the category.
\begin{center} 
	\includegraphics*[width=\textwidth]{W2_2.png} 
\end{center}

[FILL friday lecture]

\section*{Human Accelerated Regions (HARS)}
What regions are highly constrained across vertebrates, but show differences in the human lineage (usually relaive to chimpanzee)?
\begin{itemize}
	\item Reporter Assay: what patterns of gene expression are driven by a given sequence (draw)?
	\item HARs tend to be developmental enhancers.
	\item High Throughput Reporer Assay: testing thousands of sequences in cultured neurons reveals differences between human and chimpanzee sequence.  Most differences lead to higher neural activity (enhancers) in humans.
\end{itemize}

\section*{Linking Genotype to Phenotype}
Can we relate the rate of evolution ``correlate'' with enhancer activity?
\begin{itemize}
	\item Can we use linear regression to compare two continuous variables across species (compare body size with evolutionary rate across mammals?)
	\item Phylogenetic regression controls for evolutionary relationships
\end{itemize}

\subsection*{Relative Evolutionary Rates - RER Coverage}
Does the rate of evolution ``correlate'' with enhancer activity?
\begin{enumerate}
	\item Align a set of sequences together (amino acid or nucleotide)
	\item Calculate the rate of evolution along each branch in the tree
	\item Compare the branches in which a behavior evolved to other branches
	\item Check for a significant association (faster or slower) evolution along the highlighted branches
\end{enumerate}

[Fill]

\subsection*{Example: Vocal Learning}
Vocal learning is the ability to modify acoustic output based on auditory experience, enabling to imitate and produce new sounds.
\begin{itemize}
	\item Vocal learning has undergone convergent evolution, having evolved independently in multiple lineages of mammals and birds.
\end{itemize}

\subsubsection*{Modules of Vocal Learning}
If an animal falls under all three modules, it is considered a vocal learner.
\begin{itemize}
	\item \textbf{Vocal Versatility}
	\begin{itemize}
	    \item Acoustic repertoire complexity
	    \item Source-filter model of production
	    \item Pitch shifting / pitch matching
	    \item Vocal imitation / mimicry
	    \item Improvisation / innovation
    \end{itemize}
	\item \textbf{Vocal Production Variability}
	\begin{itemize}
	    \item Ontogenetic / developental variability
	    \item Babbling / subsong / plastic song
	    \item Critical periods for vocal plasticity
	    \item Closed- and open-ended learning
    \end{itemize}
	\item \textbf{Vocal Coordination}
	\begin{itemize}
	    \item Antiphonal calling
	    \item Jamming avoidance
	    \item Chorusing
	    \item Turn-taking
	    \item Duetting
	    \item Rhythmic entrainment
    \end{itemize}
\end{itemize}

\subsubsection*{Molecular Brain Similarities}
[Fill image]

\subsubsection*{Vocal Learning Enhancer Convergent Evolution}
\begin{itemize}
	\item Most enhancers shows lower predicted activity in vocal learning species.
	\item Enhancers show a tendency to be near autism-associated genes.
\end{itemize}


\end{document}